\section{Estructuras algebraicas fundamentales y categorías}

\textbf{Referencia principal:} \texttt{First Semester/Algebra/ref/doc.pdf}

\textbf{Capítulo:} 1

\textbf{Secciones:} 1.1 y 1.2

Este set está alineado con las definiciones y ejemplos del capítulo: leyes de composición, magma, semigrupo, monoide, grupo y nociones básicas de categoría.

\subsection*{Parte A --- Leyes de composición y estructuras algebraicas}

\begin{enumerate}[label=\textbf{\arabic*.}]
    \item Sea \( A = \mathbb{Z} \). Defina la operación
    \[
        a * b = a + b + 1.
    \]
    Determine si \( (A, *) \) es magma, semigrupo, monoide o grupo. Si existe elemento identidad, encuéntrelo. Si es grupo, halle el inverso de un elemento arbitrario.

    \item Considere los siguientes ejemplos y clasifique cada uno como magma, semigrupo, monoide o grupo. Indique además si la operación es conmutativa:
    \begin{enumerate}[label=\alph*)]
        \item \( (\mathbb{N}, +) \)
        \item \( (\mathbb{Z}, +) \)
        \item \( (\mathbb{N}, \cdot) \)
        \item \( (\mathcal{P}(X), \cup) \)
        \item \( (\mathcal{P}(X), \cap) \)
        \item \( (\mathrm{End}(X), \circ) \), donde \( \mathrm{End}(X) \) es el conjunto de funciones \( X \to X \)
    \end{enumerate}

    \item Dé un ejemplo de:
    \begin{enumerate}[label=\alph*)]
        \item una ley de composición entre dos conjuntos distintos que NO sea operación interna;
        \item una operación interna que no sea asociativa;
        \item un semigrupo que no sea monoide;
        \item un monoide que no sea grupo.
    \end{enumerate}

    \item En el conjunto \( A = \{0,1,2\} \), defina la operación
    \[
        a \star b = \max\{a,b\}.
    \]
    Verifique si \( (A,\star) \) es un monoide. Determine su elemento identidad y diga si es grupo.

    \item En el conjunto \( A = \{0,1,2\} \), defina la operación
    \[
        a \diamond b = \min\{a+b,2\}.
    \]
    Estudie si la operación es asociativa. Si lo es, determine si \( (A,\diamond) \) es monoide o grupo.

    \item Pruebe que en un monoide el elemento identidad es único.

    \item Pruebe que en un grupo el inverso de cada elemento es único.

    \item Sea \( G \) un grupo. Demuestre que para todo \( a \in G \),
    \[
        (a^{-1})^{-1} = a
        \qquad \text{y} \qquad
        (ab)^{-1} = b^{-1}a^{-1}.
    \]

    \item Demuestre que \( (\mathbb{Z}_n, +) \) es un grupo abeliano para todo \( n \geq 1 \).

    \item Usando la convención del texto, considere
    \[
        (\mathbb{Z}_n)^* = \{[a] \in \mathbb{Z}_n : [a] \neq [0]\}
    \]
    con la multiplicación módulo \( n \).
    \begin{enumerate}[label=\alph*)]
        \item Estudie los casos \( n=4 \), \( n=6 \) y \( n=7 \).
        \item Determine en cuáles casos se obtiene un grupo.
        \item Explique qué relación observa con la primalidad de \( n \).
    \end{enumerate}

    \item Dé un ejemplo de operación externa \( A \times B \to B \) y explique por qué NO puede considerarse una operación interna en \( B \).

    \item Sea \( M \) un semigrupo que tiene identidad por la izquierda \( e \) y tal que cada elemento \( a \in M \) tiene inverso por la izquierda \( b \in M \), es decir, \( ba = e \). Demuestre que \( M \) es un grupo.
\end{enumerate}

\subsection*{Parte B --- Introducción a categorías}

\begin{enumerate}[label=\textbf{\arabic*.}, start=13]
    \item Escriba con precisión la definición de categoría identificando:
    \begin{enumerate}[label=\alph*)]
        \item objetos,
        \item morfismos,
        \item composición,
        \item identidades,
        \item asociatividad.
    \end{enumerate}

    \item En la categoría \textbf{Set}, considere
    \[
        A=\{1,2\}, \qquad B=\{a,b,c\}, \qquad C=\{\alpha,\beta\}.
    \]
    Defina funciones \( f:A\to B \) y \( g:B\to C \), calcule \( g\circ f \) y verifique explícitamente la propiedad de identidad usando \( id_A \) e \( id_B \).

    \item Demuestre que el morfismo identidad de un objeto en una categoría es único.

    \item Sea \( f:A\to B \) un morfismo en una categoría. Suponga que existe \( g:B\to A \) con
    \[
        g\circ f = id_A
    \]
    y \( h:B\to A \) con
    \[
        f\circ h = id_B.
    \]
    Demuestre que \( g=h \).

    \item En la categoría \textbf{Set}, pruebe que un morfismo es un isomorfismo si y solo si es biyectivo.

    \item En la categoría \textbf{Mag}, tome
    \[
        (\mathbb{Z},+) \qquad \text{y} \qquad (\mathbb{Z},+).
    \]
    Estudie si las aplicaciones siguientes son homomorfismos de magmas:
    \begin{enumerate}[label=\alph*)]
        \item \( f(x)=2x \)
        \item \( g(x)=x+1 \)
        \item \( h(x)=-x \)
    \end{enumerate}

    \item En la categoría \( \mathbf{Vect}_K \), explique por qué la composición de aplicaciones lineales vuelve a ser lineal. Luego verifique esa propiedad para dos transformaciones lineales concretas entre espacios \( K^2 \) y \( K^3 \).

    \item Dé un ejemplo de un morfismo en \textbf{Top} que sea continuo pero no sea isomorfismo. Explique con cuidado por qué falla la condición de isomorfismo.
\end{enumerate}

\subsection*{Parte C --- Conexión entre monoides y categorías}

\begin{enumerate}[label=\textbf{\arabic*.}, start=21]
    \item Sea \( (M,*) \) un monoide. Construya una categoría \( \mathcal{C}_M \) con un solo objeto \( \bullet \) tal que
    \[
        \mathrm{Hom}_{\mathcal{C}_M}(\bullet,\bullet)=M.
    \]
    Defina la composición usando la operación de \( M \) y verifique los axiomas de categoría.

    \item Recíprocamente, sea \( \mathcal{C} \) una categoría con un único objeto \( A \). Demuestre que
    \[
        \mathrm{Hom}_{\mathcal{C}}(A,A)
    \]
    tiene estructura de monoide con la composición.

    \item Explique por qué un grupo puede verse como una categoría con un solo objeto en la que todos los morfismos son isomorfismos.

    \item Compare las siguientes tres ideas y explique sus diferencias con claridad:
    \begin{enumerate}[label=\alph*)]
        \item igualdad de objetos;
        \item isomorfismo de objetos;
        \item igualdad de morfismos.
    \end{enumerate}
\end{enumerate}

\subsection*{Reto final}

\begin{enumerate}[label=\textbf{\arabic*.}, start=25]
    \item Redacte un cuadro comparativo donde aparezcan las estructuras
    \[
        \text{magma} \to \text{semigrupo} \to \text{monoide} \to \text{grupo}
    \]
    indicando en cada paso qué propiedad adicional se exige. Luego explique cómo la noción de categoría abstrae la composición y la identidad que ya aparecen en estas estructuras.
\end{enumerate}

% ------------------------------------------------------------------
\subsection*{Parte D --- Problemas resueltos}
% ------------------------------------------------------------------

\begin{enumerate}[label=\textbf{\arabic*.}, start=26]

    % ---------------------------------------------------------------
    \item Demuestre que \((\mathbb{Q}^{*}, \cdot)\), donde
    \(\mathbb{Q}^{*} := \mathbb{Q} \setminus \{0\}\), es un grupo abeliano.

    \textbf{Solución.}
    \begin{itemize}
        \item \textit{Clausura.} Si \(p, q \in \mathbb{Q}^{*}\), entonces \(pq \in \mathbb{Q}\) y \(pq \neq 0\) (pues \(\mathbb{Q}\) no tiene divisores de cero). Luego \(pq \in \mathbb{Q}^{*}\).
        \item \textit{Asociatividad.} La multiplicación en \(\mathbb{Q}\) es asociativa.
        \item \textit{Identidad.} El número \(1 \in \mathbb{Q}^{*}\) satisface \(1 \cdot q = q \cdot 1 = q\) para todo \(q \in \mathbb{Q}^{*}\).
        \item \textit{Inversos.} Para cada \(q \in \mathbb{Q}^{*}\), el número \(q^{-1} = 1/q \in \mathbb{Q}^{*}\) satisface \(q \cdot q^{-1} = 1\).
        \item \textit{Conmutatividad.} \(pq = qp\) para todo \(p,q \in \mathbb{Q}^{*}\).
    \end{itemize}
    Por tanto \((\mathbb{Q}^{*}, \cdot)\) es un grupo abeliano. \hfill\(\square\)

    % ---------------------------------------------------------------
    \item Demuestre que \(GL_n(\mathbb{R})\), el conjunto de matrices invertibles \(n \times n\) con coeficientes reales, es un grupo bajo la multiplicación de matrices.

    \textbf{Solución.}
    \begin{itemize}
        \item \textit{Clausura.} Si \(A, B \in GL_n(\mathbb{R})\), entonces \(\det(AB) = \det A \cdot \det B \neq 0\), luego \(AB \in GL_n(\mathbb{R})\).
        \item \textit{Asociatividad.} La multiplicación de matrices es asociativa.
        \item \textit{Identidad.} La matriz identidad \(I_n\) satisface \(I_n A = A I_n = A\) y \(\det(I_n)=1 \neq 0\).
        \item \textit{Inversos.} Por definición, cada \(A \in GL_n(\mathbb{R})\) admite \(A^{-1}\) con \(\det(A^{-1}) = (\det A)^{-1} \neq 0\).
    \end{itemize}
    Por tanto \(GL_n(\mathbb{R})\) es un grupo. (No es abeliano para \(n \geq 2\) en general.) \hfill\(\square\)

    % ---------------------------------------------------------------
    \item Sea \(f : (G, *) \to (H, \cdot)\) un homomorfismo de grupos. Demuestre que:
    \begin{enumerate}[label=\alph*)]
        \item \(f(e_G) = e_H\),
        \item \(f(a^{-1}) = f(a)^{-1}\) para todo \(a \in G\).
    \end{enumerate}

    \textbf{Solución.}
    \begin{enumerate}[label=\alph*)]
        \item Como \(f\) es homomorfismo:
        \[
            f(e_G) = f(e_G * e_G) = f(e_G) \cdot f(e_G).
        \]
        Multiplicando ambos lados a la izquierda por \(f(e_G)^{-1}\):
        \[
            e_H = f(e_G)^{-1} \cdot f(e_G) \cdot f(e_G) = e_H \cdot f(e_G) = f(e_G).
        \]
        \item Para \(a \in G\):
        \[
            f(a) \cdot f(a^{-1}) = f(a * a^{-1}) = f(e_G) = e_H.
        \]
        Análogamente \(f(a^{-1}) \cdot f(a) = e_H\). Por unicidad del inverso, \(f(a^{-1}) = f(a)^{-1}\). \hfill\(\square\)
    \end{enumerate}

    % ---------------------------------------------------------------
    \item Considere \(A = \{e, a, b, c\}\) con la siguiente tabla de Cayley:
    \[
    \begin{array}{c|cccc}
        \cdot & e & a & b & c \\ \hline
        e & e & a & b & c \\
        a & a & e & c & b \\
        b & b & c & e & a \\
        c & c & b & a & e
    \end{array}
    \]
    \begin{enumerate}[label=\alph*)]
        \item Verifique que \((A, \cdot)\) es un grupo.
        \item Determine si es abeliano.
        \item Identifique el inverso de cada elemento.
    \end{enumerate}

    \textbf{Solución.}
    \begin{enumerate}[label=\alph*)]
        \item \textit{Clausura}: todos los productos están en \(A\) (la tabla es cerrada). \textit{Asociatividad}: se puede verificar directamente (es \(\mathbb{Z}_2 \times \mathbb{Z}_2\), conocido como grupo de Klein). \textit{Identidad}: \(e\) actúa como neutro. \textit{Inversos}: ver apartado c).
        \item La tabla es simétrica respecto a la diagonal, luego \(x \cdot y = y \cdot x\) para todos \(x,y \in A\). El grupo \textbf{es abeliano}.
        \item De la tabla: \(e^{-1}=e\), \(a^{-1}=a\), \(b^{-1}=b\), \(c^{-1}=c\). Todo elemento es su propio inverso. \hfill\(\square\)
    \end{enumerate}

    % ---------------------------------------------------------------
    \item Construya explícitamente la categoría \(\mathcal{C}_{(\mathbb{Z},+)}\) asociada al monoide \((\mathbb{Z}, +)\) con un único objeto \(\bullet\), y verifique los axiomas de categoría.

    \textbf{Solución.}
    \begin{itemize}
        \item \textit{Objeto}: \(\mathrm{Ob}(\mathcal{C}) = \{\bullet\}\).
        \item \textit{Morfismos}: \(\mathrm{Hom}(\bullet, \bullet) = \mathbb{Z}\). Cada entero \(n \in \mathbb{Z}\) es un morfismo \(n : \bullet \to \bullet\).
        \item \textit{Composición}: dados \(m, n \in \mathbb{Z}\), definimos \(n \circ m := m + n\).
        \item \textit{Identidad}: el morfismo identidad de \(\bullet\) es \(0 \in \mathbb{Z}\), pues \(0 + n = n + 0 = n\) para todo \(n\).
        \item \textit{Asociatividad}: \((p \circ q) \circ r = (q+p)+r = q + p + r\) y \(p \circ (q \circ r) = p + (r+q) = p+q+r\). Como la suma en \(\mathbb{Z}\) es asociativa, se cumple la asociatividad.
    \end{itemize}
    Así \(\mathcal{C}_{(\mathbb{Z},+)}\) es una categoría. Además, como todo \(n \in \mathbb{Z}\) tiene inverso \(-n\), todos los morfismos son isomorfismos, lo cual refleja que \((\mathbb{Z},+)\) es un grupo. \hfill\(\square\)

    % ---------------------------------------------------------------
    \item Sea \((M, *)\) un monoide conmutativo. Defina la relación \(\sim\) sobre \(M \times M\) por
    \[
        (a, b) \sim (c, d) \iff \exists\, t \in M : a * d * t = b * c * t.
    \]
    \begin{enumerate}[label=\alph*)]
        \item Demuestre que \(\sim\) es una relación de equivalencia.
        \item Explique por qué esta construcción es el primer paso de la construcción de Grothendieck para obtener un grupo a partir de \((M, *)\).
    \end{enumerate}

    \textbf{Solución.}
    \begin{enumerate}[label=\alph*)]
        \item
        \begin{itemize}
            \item \textit{Reflexividad}: \((a,b) \sim (a,b)\) pues \(a*b*e = b*a*e\) (usando conmutatividad, con \(t = e\)).
            \item \textit{Simetría}: si \((a,b) \sim (c,d)\) con testigo \(t\), entonces \(a*d*t = b*c*t\), que reescrito es \(c*b*t = d*a*t\), así \((c,d) \sim (a,b)\) con el mismo testigo.
            \item \textit{Transitividad}: si \((a,b) \sim (c,d)\) con testigo \(t_1\) y \((c,d) \sim (p,q)\) con testigo \(t_2\), entonces
            \[
                a*d*t_1 = b*c*t_1, \quad c*q*t_2 = d*p*t_2.
            \]
            Multiplicando convenientemente (y usando conmutatividad) se obtiene \((a,b) \sim (p,q)\) con testigo \(t_1 * t_2 * d\).
        \end{itemize}
        \item Las clases de equivalencia \([a,b]\) representan formalmente la ``diferencia'' \(a - b\) (o cociente \(a/b\) en el caso multiplicativo). Dotando al cociente \((M \times M)/{\sim}\) de la operación \([a,b] * [c,d] := [a*c, b*d]\), se obtiene un grupo: el \textbf{grupo de Grothendieck} de \(M\). Para \(M = (\mathbb{N},+)\) esta construcción produce \(\mathbb{Z}\). \hfill\(\square\)
    \end{enumerate}

\end{enumerate}
